% Part: incompleteness
% Chapter: representability-in-q
% Section: introduction

\documentclass[../../../include/open-logic-section]{subfiles}

\begin{document}

\olfileid[id]{inc}{req}{int}
\olsection{Pengantar}

Teorema-teorema ketaklengkapan berlaku bagi teori-teori yang memungkinkan
fakta-fakta dasar tentang fungsi yang dapat dihitung dinyatakan dan
dibuktikan. Kita akan menguraikan teori semacam itu yang sangat minimal,
bernama ``$\Th{Q}$'' (atau kadang-kadang ``$Q$ Robinson,'' menurut Raphael
Robinson). Kita akan menjelaskan apa artinya suatu fungsi
\emph{dapat direpresentasikan} dalam $\Th{Q}$, lalu membuktikan hal berikut:
\begin{quote}
  Suatu fungsi dapat direpresentasikan dalam $\Th{Q}$ jika dan hanya jika
  fungsi itu dapat dihitung.
\end{quote}
Antara lain, hal ini memberi kita model komputabilitas lain. Namun, kita juga
akan menggunakannya untuk menunjukkan bahwa himpunan
$\Setabs{!A}{\Th{Q} \Proves !A}$ tidak dapat diputuskan, dengan mereduksi
masalah penghentian kepadanya. Pada akhir pembahasan, kita bahkan akan telah
membuktikan hal-hal yang jauh lebih kuat daripada ini.

Bahasa $\Th{Q}$ ialah bahasa aritmetika; $\Th{Q}$ terdiri atas
aksioma-aksioma berikut (yang digunakan bersama aksioma dan aturan logika orde
pertama lainnya dengan !!{identity}):
\begin{align*}
& \lforall[x][\lforall[y][(\eq[x'][y'] \lif \eq[x][y])]] \tag{$!Q_1$}\\
& \lforall[x][\eq/[\Obj 0][x']] \tag{$!Q_2$}\\
& \lforall[x][(\eq[x][\Obj 0] \lor \lexists[y][\eq[x][y']])] \tag{$!Q_3$}\\
& \lforall[x][\eq[(x + \Obj 0)][x]] \tag{$!Q_4$}\\
& \lforall[x][\lforall[y][\eq[(x + y')][(x + y)']]] \tag{$!Q_5$}\\
& \lforall[x][\eq[(x \times \Obj 0)][\Obj 0]] \tag{$!Q_6$}\\
& \lforall[x][\lforall[y][\eq[(x \times y')][((x \times y) + x)]]] \tag{$!Q_7$}\\
& \lforall[x][\lforall[y][(x < y \liff \lexists[z][\eq[(z' + x)][y]])]] \tag{$!Q_8$}
\end{align*}
Untuk setiap bilangan asli $n$, definisikan numeral $\num{n}$ sebagai suku
$\Obj{0}^{\prime\prime\ldots\prime}$ dengan jumlah tanda prima seluruhnya
sebanyak $n$. Jadi, $\num{0}$ ialah !!{constant}~$\Obj{0}$ itu sendiri,
$\num{1}$ ialah $\Obj{0}'$, $\num{2}$ ialah $\Obj{0}''$, dan seterusnya.

Sebagai teori aritmetika, $\Th{Q}$ \emph{sangat} lemah; misalnya, kita bahkan
tidak dapat membuktikan fakta sangat sederhana seperti
$\lforall[x][\eq/[x][x']]$ atau $\lforall[x][\lforall[y][(x + y) = (y +
    x)]]$. Namun, akan kita lihat bahwa sebagian besar alasan mengapa
$\Th{Q}$ begitu menarik ialah \emph{karena} teori itu begitu lemah. Bahkan,
teori itu persis cukup kuat agar teorema ketaklengkapan berlaku. Alasan lain
$\Th{Q}$ menarik ialah karena teori itu memiliki himpunan aksioma yang
\emph{berhingga}.

Teori yang lebih kuat daripada $\Th{Q}$ (disebut \emph{aritmetika Peano}
$\Th{PA}$) diperoleh dengan menambahkan skema induksi pada~$\Th{Q}$:
\[
(!A(\Obj 0) \land \lforall[x][(!A(x) \lif !A(x'))]) \lif \lforall[x][!A(x)]
\]
dengan $!A(x)$ sebarang formula. Jika $!A(x)$ memuat !!{variable}s bebas
selain $x$, kita tambahkan kuantor universal di bagian depan untuk mengikat
semuanya (sehingga contoh skema induksi yang bersesuaian merupakan
!!a{sentence}). Misalnya, jika $!A(x, y)$ juga memuat !!{variable}~$y$ yang
bebas, contoh yang bersesuaian ialah
\[
\lforall[y][((!A(\Obj 0) \land \lforall[x][(!A(x) \lif !A(x'))]) \lif
  \lforall[x][!A(x)])]
\]
Dengan menggunakan contoh-contoh skema induksi, kita dapat membuktikan jauh
lebih banyak dari aksioma-aksioma~$\Th{PA}$ daripada dari aksioma-aksioma
$\Th{Q}$. Bahkan, diperlukan cukup banyak kerja untuk menemukan pernyataan
``alami'' tentang bilangan asli yang tidak dapat dibuktikan dalam aritmetika
Peano!{}

\begin{defn}
\ollabel{defn:representable-fn}
  Fungsi $f(x_0,\ldots,x_k)$ dari bilangan asli ke bilangan asli dikatakan
  {\em dapat direpresentasikan dalam $\Th{Q}$} jika terdapat formula
  $!A_f(x_0,\dots,x_k,y)$ sedemikian sehingga setiap kali
  $f(n_0,\dots,n_k) = m$, $\Th{Q}$ membuktikan
\begin{enumerate}
\item\ollabel{defn:rep:a} $!A_f(\num{n_0}, \dots, \num{n_k}, \num{m})$
\item\ollabel{defn:rep:b} $\lforall[y][(!A_f(\num{n_0}, \dots,
\num{n_k}, y) \lif \num{m} = y)]$.
\end{enumerate}
\end{defn}

Ada cara lain untuk menyatakan definisi ini; misalnya, secara ekuivalen kita
dapat mensyaratkan bahwa $\Th{Q}$ membuktikan
$\lforall[y][(!A_f(\num{n_0}, \dots, \num{n_k}, y) \liff
\eq[y][\num{m}])]$.

\begin{thm}
\ollabel{thm:representable-iff-comp}
Suatu fungsi dapat direpresentasikan dalam $\Th{Q}$ jika dan hanya jika
fungsi itu dapat dihitung.
\end{thm}

Ada dua arah dalam membuktikan teorema tersebut. Arah dari kiri ke kanan
cukup langsung setelah aritmetisasi sintaksis tersedia. Arah lainnya
memerlukan lebih banyak kerja. Gagasan dasarnya sebagai berikut: kita memilih
``rekursif umum'' sebagai cara untuk membuat istilah ``dapat dihitung'' tepat,
lalu menunjukkan bahwa setiap fungsi rekursif umum dapat direpresentasikan
dalam~$\Th{Q}$.
Ingat bahwa suatu fungsi bersifat rekursif umum jika dapat didefinisikan dari
$\Zero$, fungsi penerus~$\Succ$, dan fungsi-fungsi proyeksi~$\Proj{n}{i}$,
dengan menggunakan komposisi, rekursi primitif, dan pencarian tak berbatas
pada fungsi reguler. Jadi, salah satu cara menunjukkan bahwa setiap fungsi
rekursif umum dapat direpresentasikan dalam~$\Th{Q}$ ialah menunjukkan bahwa
fungsi-fungsi dasar dapat direpresentasikan, dan bahwa setiap kali beberapa
fungsi dapat direpresentasikan, fungsi-fungsi yang didefinisikan darinya
dengan komposisi, rekursi primitif, dan pencarian tak berbatas pada fungsi
reguler juga dapat direpresentasikan. Dengan kata lain, kita dapat menunjukkan
bahwa fungsi-fungsi dasar dapat direpresentasikan dan bahwa fungsi-fungsi yang
dapat direpresentasikan ``tertutup terhadap'' komposisi, rekursi primitif,
dan pencarian tak berbatas pada fungsi reguler. Hal ini menjamin bahwa setiap
fungsi rekursif umum dapat direpresentasikan.

Ternyata langkah untuk menunjukkan bahwa fungsi-fungsi yang dapat
direpresentasikan tertutup
terhadap rekursi primitif itu sulit. Untuk menghindari langkah ini, pertama
kita menunjukkan bahwa sebenarnya rekursi primitif dapat ditiadakan. Artinya,
kita menunjukkan bahwa setiap fungsi rekursif umum dapat didefinisikan dari
fungsi-fungsi dasar hanya dengan komposisi dan pencarian tak berbatas pada
fungsi reguler. Untuk itu,
kita menunjukkan bahwa rekursi primitif sesungguhnya dapat dilakukan dengan
suatu pencarian tak berbatas tertentu pada fungsi reguler. Akan tetapi, agar
hal ini berhasil, kita
harus menambahkan beberapa fungsi dasar lain: penjumlahan, perkalian, dan
fungsi karakteristik relasi identitas~$\Char{=}$. Kemudian, kita dapat
membuktikan teorema tersebut dengan menunjukkan bahwa semua fungsi dasar
\emph{ini} dapat direpresentasikan dalam~$\Th{Q}$ dan bahwa fungsi-fungsi yang
dapat direpresentasikan tertutup terhadap komposisi dan pencarian tak berbatas
pada fungsi reguler.

\end{document}
